%%% main.tex --- Hyperion: HTMX-First, Node-Free Full-Stack Web in Common Lisp
%%%
%%% Portable skeleton: compiles with a stock TeX Live (base `article` + common
%%% packages). Build:  latexmk -pdf main.tex   (or `make`).

\documentclass[11pt,a4paper]{article}

\usepackage[utf8]{inputenc}
\usepackage[T1]{fontenc}
\usepackage{lmodern}
\usepackage[margin=1in]{geometry}
\usepackage{microtype}
\usepackage{booktabs}
\usepackage{listings}
\usepackage{xcolor}
\usepackage[numbers,sort&compress]{natbib}
\usepackage[hidelinks]{hyperref}

\lstset{basicstyle=\ttfamily\small, columns=fullflexible, keepspaces=true,
  showstringspaces=false, frame=single, framesep=4pt, breaklines=true}

\title{\textbf{Hyperion}: HTMX-First, Node-Free\\ Full-Stack Web in Common Lisp}
\author{Bob\\ \small\texttt{eternal.recursion@proton.me}}
\date{\today}

\begin{document}
\maketitle

\begin{abstract}
Modern web development is dominated by a JavaScript build ecosystem whose
complexity is largely incidental to the problem of delivering interactive
hypermedia. \textsc{HTMX} demonstrates that a server rendering HTML fragments can
recover much of that interactivity without a single-page-application toolchain.
We argue that Common Lisp is an unusually strong host for this model---its macros
afford an ergonomic HTML DSL, its condition system makes failure legible, and its
\emph{live image} enables editing a running program---and present \textsc{Hyperion},
a full-stack web framework biased toward ``Common Lisp all the way down.''
Hyperion pairs a compile-time-checked HTMX vocabulary (in Coalton) with an open,
effectful rendering layer (CLOS + Spinneret), authors all client JavaScript via
Parenscript (no Node.js), generates CSS from a Lisp DSL, and provides a
Figwheel-style hot-reload development loop. We situate Hyperion within a small
ecosystem whose thesis is to make Coalton a first-class ``modern Lisp.''
\end{abstract}

\section{Introduction}
\label{sec:intro}
% The JS tooling tax; HTMX as a corrective; why a Lisp; the "CL all the way down" aim.

\section{The JavaScript Tooling Tax}
\label{sec:tax}
% Build chains, transpilers, bundlers; incidental vs essential complexity.

\section{The HTMX Case}
\label{sec:htmx}
% Hypermedia + fragment rendering; how far it reaches toward SPA interactivity.

\section{Typed HTMX: a Coalton Core and a CLOS Shell}
\label{sec:typed}
% ADTs for swaps/triggers/targets/verbs; a Duration type; static Coalton dispatch
% vs open CLOS rendering; generic components, never domain ones.

\section{Live Development}
\label{sec:live}
% The hot-reload loop: recompile into the running image, rebuild the server around
% preserved state, browser refreshes; compile errors surfaced in-page; the contrast
% with restart-and-replay stacks.

\section{Common Lisp to JavaScript, without Node}
\label{sec:cljs}
% Parenscript today; a typed Coalton->TS path as a far horizon.

\section{The Ecosystem}
\label{sec:ecosystem}
% cons (tooling), hyperion (web), praxeon (agents), a Coalton-first FP library;
% the shared house style; Coalton as a first-class modern Lisp.

\section{Evaluation Plan}
\label{sec:eval}
% Expressiveness vs a Node/React equivalent; hot-reload latency; component coverage.

\section{Conclusion}
\label{sec:conclusion}

\bibliographystyle{plainnat}
\bibliography{references}

\end{document}
